\chapter{Contesto e motivazione}
\captionsetup[figure]{font=small,labelfont=bf}
\label{app:contesto}

\section{Contesto esteso dell'introduzione}
\label{app:introduzione}
\begin{onehalfspace}

Questa sezione approfondisce il contesto del Capitolo~\ref{chap:introduzione}: l'origine del paradigma quantistico, la vulnerabilità di RSA all'algoritmo di Shor, la risposta della crittografia post-quantistica e il divario fra i dispositivi attuali e il calcolo fault-tolerant.

\subsection{Dal Limite Classico al Paradigma Quantistico}
\label{app:intro:origini}
\label{app:intro:algoritmi}

La crescita esponenziale del numero di transistor per chip, descritta nel 1965 da Gordon Moore come previsione industriale, si avvicina a limiti fisici: i transistor moderni misurano pochi nanometri, una scala alla quale gli effetti quantistici diventano inevitabili. A questo limite se ne affianca uno matematico: per alcuni problemi, come la fattorizzazione di grandi interi, nessun aumento di potenza di calcolo compensa algoritmi la cui richiesta di risorse cresce troppo rapidamente con la dimensione dell'input.

Il calcolo quantistico nasce dall'idea di usare sovrapposizione, entanglement e interferenza come risorse computazionali. Richard Feynman osservò nel 1982 che la simulazione classica di sistemi quantistici può richiedere risorse esponenziali e propose una macchina governata dalle stesse leggi dei sistemi da simulare~\cite{feynman1982}; David Deutsch formalizzò nel 1985 il primo modello di computer quantistico universale~\cite{deutsch1985}. Nel 1994 Daniel Simon individuò un problema oracle con separazione esponenziale rispetto agli algoritmi classici probabilistici, e nello stesso anno Peter Shor pubblicò algoritmi polinomiali per fattorizzazione e logaritmo discreto~\cite{shor1994}: una macchina quantistica sufficientemente grande comprometterebbe quindi RSA, Diffie--Hellman ed \ac{ECDSA}. Nel 1996 Lov Grover mostrò che la ricerca non strutturata richiede $O(\sqrt{N})$ interrogazioni anziché $O(N)$~\cite{grover1996}: uno speedup quadratico che sulla crittografia simmetrica si compensa raddoppiando la lunghezza delle chiavi, mentre per RSA non esiste un rimedio analogo.

% -----------------------------------------------
\subsection{RSA e le Implicazioni Crittografiche dell'Algoritmo di Shor}
\label{app:intro:rsa}

Il sistema \textbf{RSA}, proposto nel 1977 da Rivest, Shamir e Adleman~\cite{rsa1978}, si fonda sull'asimmetria fra moltiplicare due primi grandi $p$ e $q$, operazione elementare, e ricavarli dal prodotto $N=pq$. Si calcola $\phi(N)=(p-1)(q-1)$, si sceglie $e$ coprimo con $\phi(N)$ (di solito $e=65537$) e si ricava $d\equiv e^{-1}\pmod{\phi(N)}$. La chiave pubblica è $(N,e)$, quella privata $d$: un messaggio $m<N$ si cifra come $c=m^e\bmod N$ e si decifra come $m=c^d\bmod N$, correttezza garantita dal teorema di Eulero. Calcolare $d$ da $(N,e)$ senza conoscere la fattorizzazione di $N$ è ritenuto intrattabile.

Il miglior algoritmo classico noto, il \textit{General Number Field Sieve}~\cite{lenstra1993gnfs}, ha complessità sub-esponenziale:
\begin{equation}
    T_{\text{GNFS}}(N) = O\!\left(\exp\!\left[c\cdot(\log N)^{1/3}(\log\log N)^{2/3}\right]\right).
    \label{eq:gnfs_complexity_intro}
\end{equation}
Per $N$ a 2048 bit ciò corrisponde a circa $2^{112}$ operazioni; il record pratico, RSA-250 (829 bit), ha richiesto circa 2700 anni-CPU~\cite{rsa250_2020}. Shor porta il costo a $O(n^3)$, polinomiale: le stime più recenti indicano per RSA-2048 meno di un milione di qubit fisici e meno di una settimana di calcolo~\cite{gidney2025}, sotto ipotesi architetturali precise (Sezione~\ref{app:shor_risorse}).

Le primitive vulnerabili compaiono in TLS, nei certificati X.509, in SSH e nelle firme digitali. Il rischio \textit{``harvest now, decrypt later''} consiste nel conservare oggi dati cifrati ancora sensibili in attesa di una futura macchina capace di decifrarli. Il \ac{NIST} sottolinea che non esiste una data affidabile per la comparsa di tale macchina: le scadenze 2030--2035 delle roadmap istituzionali sono obiettivi di migrazione, non previsioni~\cite{nist_pqc_explainer2026}.

% -----------------------------------------------
\subsection{Crittografia Post-Quantistica}
\label{app:intro:pqc}

Nel 2016 il \ac{NIST} ha avviato la standardizzazione della \textbf{crittografia post-quantistica} (\ac{PQC}), cioè di algoritmi eseguibili su hardware classico e resistenti agli attacchi quantistici. Nell'agosto 2024 sono stati pubblicati i primi tre standard~\cite{nist2024pqc}: \textbf{ML-KEM} (FIPS~203) per lo scambio di chiavi e \textbf{ML-DSA} (FIPS~204) per le firme, entrambi basati su reticoli, e \textbf{SLH-DSA} (FIPS~205), firma basata solo su funzioni hash. Per questi problemi non sono noti algoritmi quantistici con speedup esponenziale. La migrazione è però lenta: molti sistemi hanno cicli di vita lunghi, il materiale crittografico è più voluminoso (800 byte per la chiave ML-KEM-512 contro circa 256 per una chiave pubblica RSA-2048) e nella fase transitoria si adottano schemi ibridi, classici e post-quantistici in parallelo.

% -----------------------------------------------
\subsection{Dall'Era NISQ al Calcolo Fault-Tolerant}
\label{app:intro:nisq}

I processori attuali sono descritti come \ac{NISQ}, termine introdotto da Preskill nel 2018~\cite{preskill2018}: il numero di qubit fisici, da solo, non li rende fault-tolerant, perché decoerenza, errori di gate, readout e connettività limitano la profondità affidabile. Un qubit non può inoltre essere copiato (teorema del no-cloning~\cite{wootters1982}) né letto senza perturbarlo, per cui la protezione dagli errori non può seguire le strategie classiche di ridondanza. Anche risultati simbolici come la dimostrazione di vantaggio computazionale di Google nel 2019~\cite{arute2019}, peraltro contestata, riguardano compiti specifici e non l'esecuzione di Shor su istanze crittografiche. La ricerca procede quindi su due direzioni complementari: la \ac{QEC}, che codifica qubit logici in molti qubit fisici al costo di un overhead elevato, e la mitigazione del rumore, che riduce gli effetti degli errori senza codifica ma con costi di campionamento.

\end{onehalfspace}

% -----------------------------------------------

\section{Dettaglio degli obiettivi e del piano sperimentale}
\label{app:obiettivi}
\begin{onehalfspace}

Obiettivi, domande di ricerca, metriche e contributi sono definiti nel Capitolo~\ref{chap:obiettivi}. Questa sezione ne riporta i dettagli operativi: le campagne con i relativi preset e il protocollo statistico del confronto appaiato.

\subsection{Campagne e Validazioni}
\label{app:obiettivi:usecase}

Il disegno separa le campagne rumorose dalle validazioni ideali:

\begin{itemize}
    \item \textbf{Campagne rumorose N15}: stesso $N=15$, stessa base $a=7$, stesso circuito e due preset, \emph{riferimento} e \emph{stress}. Sono scenari uniformi illustrativi, non calibrazioni di una \ac{QPU} reale.
    \item \textbf{Validazioni ideali N21/N35}: i moltiplicatori modulari di Beauregard sono verificati senza rumore mediante truth table, ripristino dei registri ausiliari e prove QPE ideali. Il loro costo esclude una campagna Aer rumorosa equivalente; l'assenza di tale campagna non viene interpretata come probabilità rumorosa nulla.
\end{itemize}

\subsubsection{Preset delle campagne rumorose}

La Tabella~\ref{tab:finale_3_8} riporta le sole configurazioni rumorose. I simboli $\epsilon_{1q}$ e $\epsilon_{2q}$ indicano esattamente il parametro $\lambda$ passato da Aer a \texttt{depolarizing\_error}; non sono probabilità Pauli aggregate né dati di calibrazione hardware. Il modello è uniforme su tutti i qubit.


\subsubsection{Razionale per la scelta delle istanze}
\label{subsec:razionale_istanze}

\textbf{N15 --- riferimento e stress.} La base $a=7$ ha periodo $r=4$ e consente di estrarre i fattori 3 e 5. Il moltiplicatore controllato implementa $y\mapsto7y\bmod15$ ripetendo l'operazione elementare \texttt{power} volte. Compilazione, durate dei gate e composizione dei canali sono descritte nella Sezione~\ref{app:fase1_contratto}.

\textbf{N21 e N35 --- validazione ideale.} Per $N=21$, $a=2$, e $N=35$, $a=6$, si usa l'aritmetica reversibile di Beauregard. La verifica controlla l'azione sui vettori di base validi, il ritorno a zero degli ancilla e la QPE ideale; il costo circuitale che esclude queste istanze dalla campagna rumorosa è riportato nella stessa sezione. L'esclusione dal rumore è una scelta di validità e costo, non un risultato sulla robustezza delle due istanze.

\subsection{Protocollo Statistico e Test}
\label{app:obiettivi:protocollo}

\subsubsection{Appaiamento, fallimenti e pareggi}

Ogni replica genera un unico calendario di istogrammi: M1, $M_{\mathrm{TOP4}}$ e M2 analizzano gli stessi conteggi nello stesso ordine. Un fallimento entro $M_{\max}=50$ è codificato con la sentinella $M_{\max}+1$, così rimane distinto dal successo all'ultima iterazione e non scompare dai vettori inferenziali. Il confronto usa il test dei ranghi con segno di Wilcoxon con \texttt{zero\_method='pratt'} e soglia $\alpha=0.05$: M1$>$TOP-4 e M1$>$M2 sono unilaterali, TOP-4--M2 è bilaterale. Quando sono riportati più confronti della stessa famiglia, i valori $p$ vengono corretti con Holm.

In caso di conteggi uguali, l'ordine dei candidati non dipende né dall'ordine del dizionario né dal valore numerico della bitstring: è stabilito da una priorità SHA-256 dipendente dal seed della coppia replica--iterazione. Seed base, calendari distinti per simulatore e tie-breaking, configurazione e hash degli artefatti sono registrati nel manifest.

Nelle tabelle dell'Appendice~\ref{app:fase1} il fattore $\rho$ è il rapporto fra il numero medio di iterazioni di M1 e quello della strategia confrontata: $\rho>1$ indica soltanto una media inferiore per quest'ultima, mentre la significatività è stabilita sulle differenze appaiate e non dedotta dal rapporto fra medie.

\subsubsection{Mitigazione e regole decisionali}

Il readout error è incorporato nel simulatore come matrice di confusione simmetrica parametrizzata da $p_{\mathrm{ro}}$; non viene applicata una correzione post-hoc. La sensibilità a $p_{\mathrm{ro}}$ è misurata nella campagna parametrica (Sezione~\ref{subsec:sweep_secondari}).

M1 seleziona la sola moda dell'istogramma, senza sogliatura preliminare, e applica direttamente \texttt{extract\_factors}. TOP-4 estende il tentativo ai quattro esiti più frequenti. M2 usa l'etichetta TOP-1 per decidere se accettare l'istogramma e, in caso positivo, applica TOP-4: proprio il confronto con TOP-4 privo di filtro isola l'effetto del classificatore.

\end{onehalfspace}
